\section{Outcome, baseline, and the remaining target}
\label{sec:context}

This continuation gives a larger exact reduction operation for the maintained
unknot recognizer and a concrete asymptotic separation from its preceding
group-stage policy. The new operation eliminates an entire acyclic system of
singleton relations in one shared compilation. It includes relations such as
\[
 c=a b a^{-1},\qquad d=c b^{-1} a c,
\]
which are excluded by a rule restricted to expressing one generator as a
power of one other generator. The expanded replacement words can be very
long. The computation instead preserves their common subexpressions.

The principal mathematical result is an $O(M+r)$ output-grammar bound for a
batch, where $M$ counts the nodes reachable in the original relator grammar
and $r$ is the number of live generators. This bound does not contain the
depth of the dependency graph or the length of the expanded words. The
supporting insight is a disjointness property of the paths to the chosen
singleton occurrences. It bounds the combined cost of extracting all
definitions before the global substitution is performed.

The implementation is accompanied by an independent certificate checker and
measurements against a pinned complete predecessor. It also records a failed
policy experiment. Applying the larger operation at every opportunity can
make subsequent heuristic search substantially worse. The delivered public
option therefore uses a batch when it can finish immediately at an exact
one-generator endpoint; the general algebraic operation remains available
for research and supplied witnesses.

\subsection{What was already present}

The baseline is the public ProveIt repository at commit
\begin{center}\baselinecommit.\end{center}
Its most recent changes in the relevant directory were committed on
9 October 2026 at 01:42:56 UTC.
All predecessor comparisons in this package use that commit, not a small
reimplementation of the old kernel. The inspected files include
\code{compressed_search.py}, \code{compressed_words.py},
\code{primitive_projection.py}, \code{primitive_forest.py}, both source
replayers, the corresponding tests, and the maintained synthesis sections
on primitive projections, forests, and planner accounting
\cite{ProveIt2026}.

The baseline already has several important capabilities. Group words are
stored as exact binary straight-line programs. The search can perform
singleton elimination, Whitehead transformations, and relator operations;
separate source reconstruction replays its certificate. Optional primitive
projections operate on raw words, so expensive normalization can be deferred.
The forest extension permits overlapping primitive donors when each chosen
child is a unit coordinate and their dependencies are acyclic. It shares
parent-power circuits, and the planner caches immutable word summaries while
rebuilding slot and live-generator information for each current state.
These are the starting point of the present work.

The new contribution is therefore not the first use of compression, batching,
or independent replay in this project. It removes the two-generator,
coherent-sign, pure-power restriction from a particular simultaneous
operation and supplies the size argument needed when the definitions are
general words. This comparison concerns exact singleton elimination, the
unpowered case of a primitive donor. The predecessor's separate inference
from proper powers can still use source-established torsion-freeness and
remains available under its existing hypotheses. Previous certificate
formats remain accepted.

\subsection{Claims and their exact scope}

\begin{center}
\small
\begin{tabular}{p{0.29\linewidth}p{0.61\linewidth}}
\toprule
Claim & Scope established in this article \\
\midrule
Exact simultaneous elimination & Any finite group presentation with the
specified raw singleton donors and acyclic dependencies. \\
Linear grammar growth & One shared batch, retaining every other original
relator slot; length metadata still has a nonconstant bit cost. \\
Efficient independent checking & A topologically ordered list of original
donor slots and pivot labels; no replacement-word claims are trusted. \\
Asymptotic stage improvement & An explicit family of validated stabilized
unknot diagrams, relative to the pinned forest policy. \\
Polynomial positive route & Inputs that expose a checked batch to rank one
with zero exponent in every retained relation, after polynomial preparation. \\
General quasi-polynomial recognition & Not established. Global exposure,
representation growth outside this operation, and search costs remain open. \\
\bottomrule
\end{tabular}
\end{center}

The distinction in the last two rows matters. A fast certificate checker
does not imply a fast certificate producer. A fast local reduction does not
bound how long the search takes to reach its hypotheses. An infinite easy
family can demonstrate a strict improvement in the relevant operation
model without establishing a speedup for all difficult unknot diagrams.

\subsection{Relation to the literature}

Shared substitution and compressed free-group algorithms are established
techniques. Schleimer gives a constructive treatment of compressed word
problems and applies shared image descriptions to free-group automorphisms
\cite{Schleimer2008}. Haubold and Lohrey explicitly combine Tietze generator
elimination with straight-line-program substitution and additive local size
growth in their work on compressed HNN-extension problems
\cite{HauboldLohrey2011}. Kapovich's recent Proposition~3.2 gives a quantitative
layered-substitution bound for a specified sequence from a fixed finite
endomorphism family in fixed rank \cite{Kapovich2026}. Those results motivate
the data representation used here. The present proof separately charges
variable-sized donor contexts and variable-rank dependency systems; its
claims do not follow merely by citing a fixed-family substitution theorem.

The group-theoretic equivalence behind a singleton elimination is classical
Tietze elimination. We do not claim that acyclic definitions or that
equivalence are new ideas in combinatorial group theory. The contribution
here is a precise shared-grammar bound, the ordered witness that can be
checked with small summaries, their integration into the current recognizer,
and a reproducible separation from its preceding policy. No exhaustive
priority claim is made for the standalone disjoint-path lemma.

The broader geometric target must also be stated accurately. The 2021
Oxford announcement of Lackenby's work gave $n^{O(\log n)}$, while a 2024
seminar abstract specified $2^{O((\log n)^3)}$; these are distinct bounds
within quasi-polynomial time \cite{Lackenby2021Talk,Lackenby2024Talk}.
The July 2026 preprint develops a hierarchy algorithm and certification
results. Section~9 explicitly leaves several implementation costs
unestimated and bounds iteration count using hierarchy length and surface
pattern-complexity parameters. That preprint does not itself state a
quasi-polynomial running-time theorem \cite{Lackenby2026Hierarchy}. We have not
converted that work into a complete implementation with a proved
quasi-polynomial running time. The group operation in this article is a
separate constructive advance within the existing recognition pipeline.

\subsection{The group endpoint and source provenance}

For a classical knot $K\subset S^3$, its group is infinite cyclic exactly
when $K$ is the unknot; the standard implication uses the foundational
three-manifold results of Papakyriakopoulos \cite{Papakyriakopoulos1957}.
Consequently, if a sound sequence of presentation equivalences leaves one
generator and only zero-exponent relations, it supplies an unknot
certificate. A raw word on one generator need not be the empty string:
$a^t a^{-t}$ is an immediate example. Its exponent can nevertheless be
computed exactly by addition on the compressed grammar.

The public checker reconstructs the presentation from the source planar
diagram and retains every Wirtinger relation, including the redundant one.
It does not accept an arbitrary abstract presentation as a knot certificate.
The local singleton theorem needs no torsion-freeness assumption; the
source binding is still essential for turning its final cyclic-group
conclusion into a topological verdict.
